\documentclass[11pt]{article}
\usepackage{mathblatt}
% Präambel des Setzers (werkzeuge/setzer-praeambel.tex, wortgleich); dazu \gteil (Muster T2U)
\makeatletter
\setlength{\mbpfbe}{0mm}% keine Punkte (pruefung.md § 5: nur beim Prüfstein)
\newlength{\gEin}\setlength{\gEin}{0pt}
\let\gSchrift\relax
\renewcommand{\pfkreuzzeile}[1]{\par\noindent\hangindent3.2em\hangafter1\hspace*{1.6em}$\square$\ #1\par}
\newcommand{\gkreuz}[1]{$\square$\ #1\hspace{2em}}
\newcommand{\gkreuzzeile}[2]{\par\noindent\hangindent2.2em\hangafter1\mbox{#1\,$\square$}\ #2\par}% Buchstabe, Option: eine je Zeile (neben Skizze)
% links, was man ansieht, rechts, was man tut; Spaltengrenze überall an derselben Stelle (gemeinsam.md „Satz“)
\newsavebox{\gbox}\newlength{\gLinks}
% v0.10 (Lehrer 13.10.): keine starre Halbierung – die linke Spalte ist so breit wie die Figur (+ 6 mm), die
% Antwort steht direkt daneben; ist die Figur breiter als 0,62 Zeile, steht die Antwort darunter
\newcommand{\gzwei}[2]{\par\smallskip\noindent\sbox{\gbox}{#1}%
  \setlength{\gLinks}{\dimexpr\wd\gbox+6mm\relax}%
  \ifdim\gLinks>0.62\linewidth
    \usebox{\gbox}\par\smallskip\noindent\begin{minipage}{\linewidth}#2\end{minipage}%
  \else
    \begin{minipage}[t]{\gLinks}\vspace{0pt}\usebox{\gbox}\end{minipage}%
    \begin{minipage}[t]{\dimexpr\linewidth-\gLinks\relax}\vspace{0pt}\raggedright #2\end{minipage}%
  \fi\par}
% v0.9 (12.10., pruefung.md § 5 „Marke mit B, Bedingungen rechts“): drittes Argument von pfaufg = Bedingungen
% der Prüfung (v0.10: \pfohnetr = TR durchgestrichen, „FOR“), klein grau im rechten Rand auf der Grundlinie der Nummer, Zeilen mit \\
\newcommand{\gbedrand}[2]{\ifblank{#2}{}{\rlap{\hspace*{\dimexpr#1+2mm\relax}%
  \parbox[t]{15mm}{\raggedright\scriptsize\leavevmode\color{mbgrau}\strut#2}}}}
% Teil eines Bündels (a, b …): Marke links im Rand, Bedingungen rechts im Rand, auf der Zeile des Teils
\newcommand{\gteilrand}[2]{\llap{\mbpfrandmarke{#1}\hspace*{\mbpfnr}}\gbedrand{\linewidth}{#2}}
% Teil eines Bündels: Buchstabe in eigener Spalte, Inhalt eingerückt (\gEin = Einzug, damit die Spaltengrenze von
% \gzwei an derselben Stelle bleibt)
\newenvironment{gbteil}[3]{\par\smallskip\noindent\gteilrand{#2}{#3}\makebox[\mbpfteil][l]{\textbf{#1}}%
  \begin{minipage}[t]{\dimexpr\linewidth-\mbpfteil\relax}\setlength{\gEin}{\mbpfteil}\raggedright\strut\ignorespaces}%
 {\end{minipage}\par}
% Grafik höchstens so breit wie die Spalte (Bündel nebeneinander, Koordinatensystem ohne Karo)
\newsavebox{\gpbox}
\newcommand{\gpassend}[1]{\sbox{\gpbox}{#1}\ifdim\wd\gpbox>\linewidth\typeout{GPASSEND: verkleinert \the\wd\gpbox\space auf \the\linewidth}\resizebox{\linewidth}{!}{\usebox{\gpbox}}\else\usebox{\gpbox}\fi}
% Skizzenfeld (v0.10, Lehrer 13.10.): Karo ab dem linken Rand über die Zeile, „Skizze“ klein grau oben links im Feld
% (wie \pfskizzenfeld der Vorlage); der rechte Rand bleibt den Bedingungen
\newcommand{\gskizze}[1]{\par\smallskip\noindent\begin{tikzpicture}
  \clip (0,0) rectangle ({\linewidth-1mm},{#1*10mm});
  \draw[black!16,line width=0.3pt,step=5mm] (0,0) grid ({\linewidth},{#1*10mm});
  \node[anchor=north west,inner sep=2pt,font=\scriptsize,text=mbgrau,fill=white] at (0,{#1*10mm}) {Skizze};
  \end{tikzpicture}\par}
\RenewDocumentEnvironment{pfaufg}{m m m}{%
  \begin{lrbox}{\mbaufgabebox}\begin{minipage}[t]{\linewidth}%
  \gSchrift\noindent\hspace*{\gEin}\mbpfrandmarke{#1}\gbedrand{\linewidth-\gEin}{#3}%
  \begin{minipage}[t]{\mbpfnr}\strut\textbf{#2}\end{minipage}%
  \begin{minipage}[t]{\dimexpr\linewidth-\gEin-\mbpfnr\relax}\raggedright\strut\ignorespaces}%
 {\end{minipage}%
  \end{minipage}\end{lrbox}%
  \setlength{\mbaufgabehoehe}{\dimexpr\ht\mbaufgabebox+\dp\mbaufgabebox\relax}%
  \ifdefempty{\gkopfnach}{\par\addvspace{5pt}\Needspace*{\mbaufgabehoehe}}{%
    % Stufenkopf, der mit dieser (ersten) Aufgabe zusammenbleibt: Platz für beide verlangen, dann den Kopf setzen
    % (12.10., Quadratische: ein Kopf stand allein unten, die Aufgabe mit Koordinatensystem auf der nächsten Seite)
    \par\Needspace*{\dimexpr\mbaufgabehoehe+40pt\relax}\gkopfnach\gdef\gkopfnach{}\par\addvspace{5pt}}%
  \noindent\usebox{\mbaufgabebox}\par\addvspace{5pt}}
% Stufenkopf vormerken: \gkopfvor{<Kopf>} – die nächste pfaufg setzt ihn, sobald beide zusammen Platz haben
\def\gkopfnach{}
\newcommand{\gkopfvor}[1]{\gdef\gkopfnach{#1}}
% Teilaufgabe mit Skizze links in gzwei: Buchstabe in eigener Spalte, darunter Text und Skizze
\newcommand{\gteil}[2]{\makebox[\mbpfteil][l]{\textbf{#1}}\parbox[t]{\dimexpr0.5\textwidth-\gEin-\mbpfnr-5mm-\mbpfteil\relax}{\raggedright #2}}
% Stufenkopf der Vorlage (mathblatt.sty 2026-10-12: Linie, Abstand, 13 pt, \pfjahre) ohne ihren Vorhalt von
% 14 Zeilen: der Setzer fordert vor dem Kopf 0,16 Seitenhöhe (Kopf + erste Aufgabe), mehr nicht
\setlength{\mbpfstufevorhalt}{0pt}
\makeatother
\newcommand{\karo}[1]{\par\smallskip\noindent\begin{tikzpicture}
  \clip (0,0) rectangle ({\linewidth-1mm},{#1*10mm});
  \draw[black!16,line width=0.3pt,step=5mm] (0,0) grid ({\linewidth},{#1*10mm});
  \end{tikzpicture}\par}
\newcommand{\kk}[1]{$\square$\,#1\hspace{0.9em}}
\newcommand{\linie}[1]{\,{\color{black!45}\rule[-1pt]{#1}{0.4pt}}\,}
\newcommand{\pkt}{\textperiodcentered{}}
% v0.14 (14.10.): Stufenkopf mit fünf Punkten (Option = Jahre für die Bestellkarte) oder „zuletzt 19“
\NewDocumentCommand{\gpunkte}{o m}{#2}
\newcommand{\gzuletzt}[1]{zuletzt #1}
% v0.7 (12.10.): Schülerfrage unter dem Stufenkopf; Ankreuzoption als Bild (Buchstabe und Kästchen darunter)
\newcommand{\gfrage}[1]{\par\nobreak\noindent\hspace*{\mbpfnr}{\small\itshape #1}\par\nobreak\addvspace{2pt}}
\newlength{\gbildbreite}\setlength{\gbildbreite}{0.32\linewidth}% Setzer setzt sie je Reihe (vier Bilder: 0,24)
\newcommand{\gbildoption}[3]{\begin{minipage}[t][#3][b]{\gbildbreite}\centering #2\par\smallskip\mbox{#1\,$\square$}\end{minipage}}
% Lineare (12.10.): Bild mit Buchstabe ohne Kästchen (Diagramm zuordnen); Aussage mit „wahr / falsch“ zum Ankreuzen
\newcommand{\gbildname}[3]{\begin{minipage}[t][#3][b]{\gbildbreite}\centering #2\par\smallskip\mbox{\textbf{#1}}\end{minipage}}
\newcommand{\gwahrfalsch}[2]{\par\noindent\hangindent2.2em\hangafter1\makebox[2.2em][l]{(#1)}#2\hfill\mbox{wahr $\square$\quad falsch $\square$}\par}
\begin{document}

\pfheftstil
\pflaufkopf[Abitur GK \pkt{} Prüfung]{Term und Anteil am Baum}{AS50}\pfseitevon
\pfheftkopf{Term und Anteil am Baum}{}
\pfstufenzeile{Ereignis zum Term gesucht \pkt{} unbekannter Anteil gesucht}

% ===================== Stufe: Ereignis zum Term gesucht =====================
\gkopfvor{\pfstufekopf{Ereignis zum Term gesucht}{\gpunkte[25 24]{\textcolor{black!25}{$\circ$}\,\textcolor{black!80}{$\bullet$}\,\textcolor{black!80}{$\bullet$}\,\textcolor{black!25}{$\circ$}\,\textcolor{black!25}{$\circ$}}}\special{pdf:dest (stufe-ereignis-zum-term-gesucht) [@thispage /XYZ @xpos @ypos null]}}
% kein Einstieg: keine Schülerfrage in der Gliederung (Schülerfrage, -Fälle, -Form, -Weg, -Skizze; oder Frage) – Einstieg nur von dort, nie aus der Bank (pruefung.md § 5, 12.10.)
% 2025-bebb-gk-A1.3a-m1 (Original 2025-bebb-gk-A1.3a, Schwere I, rang 74; herausgelöste Fassung: keine)
\begin{pfaufg}{GK~’25\\A1.3a}{1.}{\pfohnetr}
Ein Glücksrad hat acht gleich große Sektoren: fünfmal die 3, zweimal die 2 und einmal die 1. Es wird zweimal gedreht. Deute den Term $\left(\tfrac{2}{8}\right)^2$ im Sachzusammenhang.
\gzwei{\kreisdiagramm[1.1]{3/1,2/1,3/1,3/1,2/1,3/1,1/1,3/1}}{\karo{2}}
\fusshilfe{\mbox{1:~Wahrscheinlichkeit, bei zweimaligem Drehen zweimal die 2 zu erhalten}}
\end{pfaufg}
% 2024-bebb-gk-A1.6b-m1 (Original 2024-bebb-gk-A1.6b, Schwere II, rang 75; herausgelöste Fassung: keine)
\begin{pfaufg}{GK~’24\\A1.6b}{2.}{\pfohnetr}
Bei einer Werbeaktion erhält jedes Kind einen verpackten Ball; mit einer Wahrscheinlichkeit von 30 \% ist er glitzernd gefärbt. Beschreibe ein Zufallsexperiment und ein Ereignis, dessen Wahrscheinlichkeit der Term $\left(\tfrac{7}{10}\right)^4 + 4 \cdot \left(\tfrac{7}{10}\right)^3 \cdot \tfrac{3}{10}$ angibt.
\karo{3}
\fusshilfe{\mbox{2:~mindestens drei Bälle sind nicht glitzernd}}
\end{pfaufg}
% 2025-bebb-gk-A1.6b-m1 (Original 2025-bebb-gk-A1.6b, Schwere II, rang 76; herausgelöste Fassung: keine)
\begin{pfaufg}{GK~’25\\A1.6b}{3.}{\pfohnetr}
Eine Schwimmgruppe hat 16 Kinder; 6 von ihnen haben das Abzeichen Bronze. Gib die Bedeutung des Terms $\tfrac{\binom{6}{2} \cdot \binom{10}{3}}{\binom{16}{5}}$ im Sachzusammenhang an.
\karo{2}
\fusshilfe{\mbox{3:~Wahrscheinlichkeit, dass unter fünf zufällig ausgewählten Kindern genau zwei das Abzeichen Bronze haben}}
\end{pfaufg}

% ===================== Stufe: unbekannter Anteil gesucht =====================
\gkopfvor{\pfstufekopf{unbekannter Anteil gesucht}{\gpunkte[26 24 22]{\textcolor{black!80}{$\bullet$}\,\textcolor{black!25}{$\circ$}\,\textcolor{black!80}{$\bullet$}\,\textcolor{black!25}{$\circ$}\,\textcolor{black!80}{$\bullet$}}}\special{pdf:dest (stufe-unbekannter-anteil-gesucht) [@thispage /XYZ @xpos @ypos null]}}
% kein Einstieg: keine Schülerfrage in der Gliederung (Schülerfrage, -Fälle, -Form, -Weg, -Skizze; oder Frage) – Einstieg nur von dort, nie aus der Bank (pruefung.md § 5, 12.10.)
% 2024-bebb-gk-A1.9b-m1 (Original 2024-bebb-gk-A1.9b, Schwere II, rang 77; herausgelöste Fassung: keine)
\begin{pfaufg}{GK~’24\\A1.9b}{4.}{\pfohnetr}
Bei einer Befragung erhalten 70 \% der Jugendlichen die Frage F1 „Hast du schon einmal unentschuldigt gefehlt?“, die übrigen die Frage F2 „Hast du noch nie unentschuldigt gefehlt?“. Alle antworten wahrheitsgemäß mit Ja oder Nein. Der Anteil $p$ derer, die schon einmal gefehlt haben, ist unter den F1-Befragten so groß wie unter allen. Von 1 000 Befragten antworten 360 mit Ja. Berechne auf dieser Grundlage den Anteil $p$.
\gzwei{\baumabi{F_1/{0,7}/{\text{Ja}/{p},\text{Nein}/{1-p}},F_2/{0,3}/{\text{Ja}/{1-p},\text{Nein}/{p}}}}{\karo{4}}
\fusshilfe{\mbox{4:~0,15}}
\end{pfaufg}
% 2026-bb-gk-B4e-m1 (Original 2026-bb-gk-B4e, Schwere II, rang 78; herausgelöste Fassung: keine)
\begin{pfaufg}{GK~’26\\B4e}{5.}{}
In einem Supermarkt sammeln 80 \% der Kundschaft Treuepunkte ($T$). Unter den Sammelnden sind 70 \% weiblich ($w$), unter den übrigen ist der Anteil weiblicher Personen $a$. Der Anteil der nicht weiblichen Personen in der Kundschaft beträgt 34 \%. Berechne $a$.
\gzwei{\baumabi{T/{0,8}/{w/{0,7},\overline{w}/},\overline{T}//{w/{a},\overline{w}/}}}{\karo{4}}
\fusshilfe{\mbox{5:~0,5}}
\end{pfaufg}
% 2022-bebb-gk-B4j-m1 (Original 2022-bebb-gk-B4j, Schwere III, rang 79; herausgelöste Fassung: keine)
\begin{pfaufg}{GK~’22\\B4j}{6.}{}
In einem Paketzentrum haben 12 \% der Pakete das Ziel A und 8 \% das Ziel B; 6 \% aller Pakete sind schwer ($S$). Unter den Paketen mit Ziel A sind 10 \% schwer, unter denen mit Ziel B 2,5 \%. Untersuche, ob der Anteil schwerer Pakete unter denen mit anderem Ziel kleiner, gleich oder größer als 6 \% ist.
\karo{4}
\fusshilfe{\mbox{6:~kleiner}}
\end{pfaufg}

\pfblattende{10}{26 & 25 & 24 & 23 & 22 & 21 & 20 & 19 & 18 & 17}{B4e & A1.3a A1.6b & A1.6b A1.9b &  & B4j & B4f & B4.2g & B4.1e &  & }%
{\pfausschnitt{Vierfeldertafel \pkt{} Bedingte Wahrscheinlichkeit}{Term und Anteil am Baum}{}}
\end{document}